
%%%%%%%%%%%%%%%%%%%%%%%%%%%%%%%%%%%%%%%%
%           main body
%%%%%%%%%%%%%%%%%%%%%%%%%%%%%%%%%%%%%%%%

%-----------------------------------
%	TITLE PAGE
%-----------------------------------

\title[第十一章\\
Euclid空间上的极限与连续]{\texorpdfstring{\LARGE \bfseries 第十一章\quad Euclid 空间上的极限与连续}{第十一章 Euclid 空间上的极限与连续}} % The short title appears at the bottom of every slide, the full title is only on the title page
\author[\S 11.1 Euclid 空间基本定理] % Your institution as it will appear on the bottom of every slide, maybe shorthand to save space
{\Large \bfseries \S 11.1 Euclid 空间上的基本定理}
% \author{Singular Value Decomposition} % Your name
% \date{\today} % Date, can be changed to a custom date
\date{}

%------------------------------------------------

\begin{document}

%------------------------------------------------

\begin{frame}

\titlepage % Print the title page as the first slide

\end{frame}

%------------------------------------------------

% \begin{frame}
% \frametitle{内容提要} % Table of contents slide, comment this block out to remove it
% \tableofcontents % Throughout your presentation, if you choose to use \section{} and \subsection{} commands, these will automatically be printed on this slide as an overview of your presentation
% \end{frame}

%-----------------------------------
%	PRESENTATION SLIDES
%-----------------------------------

\section{基本定义}

%------------------------------------------------

\begin{frame}{Euclid 空间}

\begin{Def}
$(\mathbb{R}, \langle \cdot, \cdot \rangle)$ 被称作 Euclid 空间 (内积空间), 其中
\begin{gather*}
\mathbb{R}^n = \left\{ (x_1, x_2, \dots, x_n) : ~ x_1, \dots, x_n \in \mathbb{R} \right\}, \\
\langle x, y \rangle = \sum\limits_{i=1}^n x_i y_i = x y^T.
\end{gather*}
\end{Def}

\pause

设 $A$ 为 $n$ 阶正定阵, 则
$$\langle x, y \rangle_A = xAy^T$$
也给出了内积空间的结构.

\pause
\vspace{0.6em}

一般地, $\mathbb{R}^n$ 上的内积指的是一个正定对称的双线性型
$$(x, y) : \mathbb{R}^n \times \mathbb{R}^n \longrightarrow \mathbb{R}$$
$(x, y) = xGy^T,$ $G$ 正定 (对称).

\end{frame}

%------------------------------------------------

\begin{frame}{Euclid 空间例子}

\begin{align*}
M_{n\times m}(\mathbb{R}) = \{ \text{所有} n\times m \text{实矩阵} \} & \rightarrow \mathbb{R}^{n\times m} \\
\begin{pmatrix}
    a_{11} & \cdots & a_{1m} \\
    \vdots & & \vdots \\
    a_{n1} & \cdots & a_{nm}
\end{pmatrix} & \mapsto (a_{11}, \cdots, a_{1m}, a_{21}, \cdots, a_{nm})
\end{align*}

\pause

那么 $\forall A, B \in M_{n\times m}(\mathbb{R}),$ 有
$$\langle A, B \rangle = \operatorname{tr} (AB^T).$$

\end{frame}

%------------------------------------------------

\section[赋范空间]{范数与赋范线性空间}

%------------------------------------------------

\begin{frame}{范数与赋范线性空间}

内积 $\longsquiggly$ 范数 \quad $\lVert x \rVert = \sqrt{\langle x, x \rangle}$

\pause

\begin{Def}
设 $V$ 为 $\mathbb{R}$ 上线性空间, $V$ 上的一个范数指的是
$$\lVert \cdot \rVert : ~ V \longleftarrow \mathbb{R}$$
满足
\begin{enumerate}
    \item 正定性: $\lVert x \rVert \geqslant 0,$ 且 $\lVert x \rVert \geqslant 0$ 当且仅当 $x = 0;$
    \item $\forall \lambda \in \mathbb{R}, \lVert \lambda x \rVert = \lvert \lambda \rvert \lVert x \rVert;$
    \item 三角不等式: $\lVert x + y \rVert \leqslant \lVert x \rVert + \lVert y \rVert.$
\end{enumerate}
$(V, \lVert \cdot \rVert)$ 称为一个赋范线性空间.
\end{Def}

\pause

\begin{prop}
由内积诱导的映射 $x \mapsto \sqrt{\rangle x, x \rangle} = \lVert x \rVert$ 是一个范数, 并且满足平行四边形等式
$$\lVert x + y \rVert^2 + \rVert x - y \rVert^2 = 2 ( \lVert x \rVert^2 + \lVert y \rVert^2 ).$$
\end{prop}

\end{frame}

%------------------------------------------------

\begin{frame}{范数与赋范线性空间}

\begin{block}{范数的例子}
$M_{n\times m}(\mathbb{R})$ 上由 $\langle A, B \rangle = \operatorname{tr} (AB^T)$ 对应的范数为
$$\lVert A \rVert_F = \sqrt{\operatorname{tr} (AA^T)} = \sqrt{\sum\limits_{i, j} \lvert a_{ij} \rvert^2}$$
称为矩阵的 Frobenius 范数.
\end{block}

\pause

矩阵范数还有很多, 例如算子范数
\vspace{-1em}
$$
A: ~ \mathbb{R}^{(m)} \rightarrow \mathbb{R}^{(n)}, ~~ x = \begin{pmatrix}
x_1 \\ \vdots \\ x_m
\end{pmatrix} \mapsto Ax,
$$
定义$\lVert A \rVert = \sup\limits_{x\neq 0} \dfrac{\lVert Ax \rVert}{\lVert x \rVert},$
\pause
或更一般的$\sup\limits_{x\neq 0} \dfrac{\lVert Ax \rVert_q}{\lVert x \rVert_p},$ 其中
$$\lVert x \rVert_p = \left( \lvert x_1 \rvert^p + \cdots + \lvert x_n \rvert^p \right)^{1/p} ~ \text{为 $p$ 范数}.$$

\pause

{\color{red}注意}, 不是所有范数都可以由某个内积所诱导,
\pause
例如 $p$ 范数, $p\neq 2.$

\end{frame}

%------------------------------------------------

\section[度量空间]{度量空间}

%------------------------------------------------

\begin{frame}{度量空间}

范数 $\longsquiggly$ 度量 \quad $d(x, y) = \lVert x - y \rVert$

\pause
\vspace{1em}

度量满足
\begin{enumerate}
\item $d(x, y) \geqslant 0;$ 且 $d(x ,y) = 0$ 当且仅当 $x = y;$
\item $d(x, y) = d(y, x);$
\item 三角不等式 $d(x, y) \leqslant d(x, z) + d(y, z).$
\end{enumerate}

\pause
\vspace{1em}

不是所有度量都能由某个范数所诱导, 例如离散度量
$$
d(x, y) = \begin{cases}
1, & x \neq y, \\ 0, & x = y
\end{cases}
$$

\end{frame}

%------------------------------------------------

\begin{frame}{度量诱导的拓扑}

度量 $\longsquiggly$ 拓扑 \quad (定义开集、闭集)

\pause

\begin{Def}[开集]
集合 $E \subset \mathbb{R}^n$ 被称作开集, 若 $E$ 中所有的点都是{\color{red}内点}, 即 $\forall x \in E,$ {\color{red} 存在 $x$ 的邻域 $O(x, \delta), \delta > 0,$ 使得 $O(x, \delta) \subset E.$}
$$\text{邻域} ~~ O(x, \delta) = \{ y: ~ d(x, y) < \delta \}.$$
\end{Def}

\pause
\vspace{1em}

去心邻域
$$\mathring{O}(x, \delta) = \{ y: ~ 0 < d(x, y) < \delta \}.$$

\pause
\vspace{1em}

\begin{Def}[闭集]
集合 $E \subset \mathbb{R}^n$ 被称作闭集, 若它的补集 $\mathbb{R}^n \setminus E$ 是开集.
\end{Def}

\end{frame}

%------------------------------------------------

\begin{frame}{度量诱导的拓扑}

不是所有的拓扑都是由某个度量诱导的, 因为所有由度量诱导的拓扑都是 Hausdorff 的 (T2).

\pause
\vspace{1em}

T2: $\forall x \neq y,$ 存在 $x$ 的邻域 $U(x),$ $y$ 的邻域 $U(y),$ 使得 $U(x) \cap U(y) = \emptyset.$

\pause
\vspace{1em}

\begin{block}{所有由度量诱导的拓扑都是 Hausdorff 的证明}
任取 $x \neq y,$ 则 $d = d(x, y) > 0.$ 取 $x$ 的邻域 $U(x) = O(x, d/3),$ $y$ 的邻域 $U(y) = O(y, d/3).$ 假设 $U(x) \cap U(y) \neq \emptyset,$ 任取 $z \in U(x) \cap U(y),$ 那么有
$$d(x, z) < d/3, ~~ d(y, z) < d/3,$$
于是根据三角不等式有
$$d(x, y) \leqslant d(x, z) + d(y, z) < 2d/3 < d,$$
矛盾. 因此必须有$U(x) \cap U(y) = \emptyset.$
\end{block}

\end{frame}

%------------------------------------------------

\begin{frame}{其余部分类型的拓扑空间}

\begin{itemize}
\item T1 (Fréchet):

$\forall x \neq y,$ 存在 $x$ 的邻域 $U(x)$ 使得 $y \not\in U(x),$ {\color{red}并且} 存在 $y$ 的邻域 $U(y)$ 使得 $x \not \in U(y).$
\item T0 (Kolmogorov):

$\forall x \neq y,$ 存在 $x$ 的邻域 $U(x)$ 使得 $y \not\in U(x),$ {\color{red}或者} 存在 $y$ 的邻域 $U(y)$ 使得 $x \not\in U(y).$
\end{itemize}

\pause

\begin{block}{非 Hausdorff 拓扑空间的例子}
Zariski拓扑
\pause
$k^n = \operatorname{Spec} k[x_1, \dots, x_n],$ $k$为代数闭域. 更一般地, 
$$\operatorname{Spec} A = \{ A ~ \text{的所有素理想} ~ \mathfrak{p} \},$$
$A$ 为一个交换环. $\operatorname{Spec} A$ 的所有闭集为
$$V(\mathfrak{a}) = \{ \mathfrak{p} \in \operatorname{Spec} A: 
\mathfrak{a} \subset \mathfrak{p} \}, ~~ \mathfrak{a} ~\text{为理想}.$$
\pause
若素理想 $\mathfrak{p}_1 \subsetneqq \mathfrak{p}_2,$ 那么任意包含点 $\mathfrak{p}_2$ 的开集必然也包含点 $\mathfrak{p}_1.$

\pause

这是因为, 设该包含点 $\mathfrak{p}_2$ 的开集是某个 $V(\mathfrak{a})$ 的补集, 那么 $\mathfrak{a} \not\subset \mathfrak{p}_2,$ 从而必然有 $\mathfrak{a} \not\subset \mathfrak{p}_1,$ 于是该开集包含 $\mathfrak{p}_1.$

\pause

这种情况下, $\operatorname{Spec}A$ 是 T0 拓扑空间.
\end{block}

\end{frame}

%------------------------------------------------

\section{Euclid 空间进一步的定义与性质}

%------------------------------------------------

\begin{frame}{收敛与极限}

\begin{Def}
设 $\{x_k\}$ 是 $\mathbb{R}^n$ 中的点列, 若存在点 $a \in \mathbb{R}^n,$ 使得 $\forall \varepsilon > 0,$ 存在正整数 $K,$ 使得 $\forall k \in K,$ 有
$$d(x_k, a) = \lVert x_k - a \rVert < \varepsilon ~~ (\text{或}~ x_k \in O(a, \varepsilon)),$$
则称点列 $\{x_k\}$ 收敛于 $a,$ 记为 $\lim\limits_{n\to\infty} x_k = a,$ $a$ 称为点列 $\{x_k\}$ 的极限.
\end{Def}

\pause

\begin{thm}
$\lim\limits_{n\to\infty} x_k = a$ 当且仅当 $\forall 1 \leqslant i \leqslant n, \lim\limits_{n\to\infty} x_i^k = a_i,$ 其中 $x_k = (x_1^k, \dots, x_n^k).$
\end{thm}

\pause

由 $d(x_k, a) = \lVert x_k - a \rVert = \sqrt{\sum\limits_{i=1}^n (x_i^k - a_i)^2}$ 知
\vspace{-0.5em}
$$
d(x_k, a) ~~ \begin{cases}
\geqslant ~~ \lvert x_i^k - a_i \rvert, ~~ \forall 1 \leqslant i \leqslant n \\
\leqslant ~~ \sum\limits_{i=1}^n \lvert x_i^k - a_i \rvert
\end{cases}
$$

\end{frame}

%------------------------------------------------

\begin{frame}{更多点集拓扑相关的概念}

之前引入开集这一概念时, 我们介绍了内点的概念. 接下来, 我们介绍更多点集拓扑相关的概念.

\pause
\vspace{1em}

\begin{itemize}[<+->]
\item 集合 $E \subset \mathbb{R}^n$ 的内点 $x:$ ~~ $\exists O(x, \delta) \subset E.$
\item 集合 $E \subset \mathbb{R}^n$ 的外点 $x:$ ~~ $\exists O(x, \delta) \subset \mathbb{R}^n \setminus E.$
\item 集合 $E \subset \mathbb{R}^n$ 的边界点 $x:$ ~~ $\forall O(x, \delta) \cap E \neq \emptyset,$ 但 $O(x, \delta) \not\in E.$ \par
集合 $E \subset \mathbb{R}^n$ 的边界 $\partial E:$ ~~ $\{ x \in \mathbb{R}^n: ~ x ~ \text{为 $E$ 的边界点} \}.$
\item 集合 $E \subset \mathbb{R}^n$ 的孤立点 $x:$ ~~ $\exists O(x, \delta)$ 使得 $O(x, \delta) \cap E = \{x\}.$
\item 集合 $E \subset \mathbb{R}^n$ 的聚点 (极限点) $x:$ ~~ $\forall \mathring{O}(x, \delta), \mathring{O}(x, \delta) \cap E \neq \emptyset.$ \par
集合 $E \subset \mathbb{R}^n$ 的导集 $E':$ ~~ $\{ x \in \mathbb{R}^n: ~ x ~ \text{为 $E$ 的极限点} \}.$
\item 集合 $E \subset \mathbb{R}^n$ 的闭包 $\overline{E}:$ ~~ $E \cup E'.$
\end{itemize}

\end{frame}

%------------------------------------------------

\begin{frame}{聚点 (极限点)}

\begin{thm}
点 $x$ 是集合 $E$ 的聚点的充要条件是, 存在点列 $x_k \in E, x_k \neq x,$ 使得 $\lim\limits_{k\to\infty} x_k = x.$ 
\end{thm}

\pause

\begin{proof}
$\Longrightarrow:$ 由定义, 取 $\delta_1 = 1,$ 则 $\mathring{O}(x, \delta_1) \cap E \neq \emptyset,$
\end{proof}

\end{frame}

%------------------------------------------------

\section{总结}

%------------------------------------------------

\begin{frame}{度量诱导的拓扑}

\end{frame}


%------------------------------------------------

\section{总结}

%------------------------------------------------

\begin{frame}{总结}

{\larger[2] \color{red} \ding{43}} 待写

\pause
\vspace{0.8em}

{\larger[2] \color{red} \ding{43}} 待写

\pause
\vspace{0.8em}

{\larger[2] \color{red} \ding{43}} 待写

\end{frame}

%------------------------------------------------

%------------------------------------------------
% % closing page
% \begin{frame}
% \HUGE{\centerline{\bfseries {\color{gray} The} {\color{zkblue} End}}}

% \pause

% \vspace{1.2em}

% \Large{\centerline{\bfseries {\color{gray}Thank} \quad {\color{zkblue} You}}}

% \end{frame}

%----------------------------------------------------------------------------------------

\end{document}
